\chapter{Applications and more realistic Hamiltonians}
\label{chap:applications}
%=============================================================================

Every system in this book so far has been a model.  The pairing model has one
interaction and one strength; the Lipkin model has two; the Hubbard model has
an on-site $U$ and nothing else.  That was the point: with a single parameter
one can turn correlation on and off at will and watch each approximation break
where it should.  But a model interaction also hides the two things that make
real many-body calculations hard.  The matrix elements were given to us in
closed form and there were only a handful of them, and the single-particle
space was finite by construction.

This chapter removes both crutches.  The system is a \emph{quantum dot}: $N$
electrons confined to two dimensions by a parabolic trap and interacting
through the bare Coulomb repulsion.  Nothing is schematic.  The interaction is
the real one, every matrix element has to be computed, the single-particle
basis is infinite and has to be truncated, and the truncation error has to be
controlled alongside the many-body error.  Quantum dots are also real: they
are fabricated in semiconductor heterostructures, their level structure is
measured, and they are among the candidate architectures for quantum
computing.

The plan is to take the same hierarchy through this system that
Chapters~\ref{chap:hartrefock}, \ref{chap:mbpt} and~\ref{chap:cctheory} took
through the pairing model -- Hartree-Fock, second-order perturbation theory,
coupled cluster with singles and doubles, and the unitary variant -- and to
measure all of them against exact diagonalisation wherever that is affordable.
The programs are the same too: \texttt{quantumdot.py} builds the one- and
two-body matrices, and everything after that is imported unchanged from
\texttt{coupledcluster.py}.  That reuse is deliberate.  A method that works
only on the model it was developed on is of no use.

The Hamiltonian of this chapter returns in the stochastic part of the book,
where variational and diffusion Monte Carlo are applied to the same dots and
compared with the results obtained here.

%----------------------------------------------------------------------
\section{The system}
\label{sec:qdsystem}

\index{quantum dot}
We consider $N$ electrons moving in two dimensions in a harmonic trap.  In
natural units, $\hbar=m_e=e=1$, the Hamiltonian splits in the usual way into a
one-body and a two-body part,
\begin{equation}
  \boxed{\;
  \hat{H} = \hat{H}_0 + \hat{V},
  \qquad
  \hat{H}_0 = \sum_{i=1}^{N}
    \left(-\tfrac12\nabla_i^2 + \tfrac12\omega^2 r_i^2\right),
  \qquad
  \hat{V} = \sum_{i<j}^{N}\frac{1}{|\bm{r}_i-\bm{r}_j|} \;}
  \label{eq:11-hamiltonian}
\end{equation}
Lengths are measured in the oscillator length
$a_0^{*}=\sqrt{\hbar/(m_e\omega)}$.  The single parameter of the problem is
the trap frequency $\omega$, and it does more work than it looks.  The kinetic
and confining energies scale as $\omega$ while the Coulomb energy scales as
$\sqrt{\omega}$, so a shallow trap is a strongly correlated system and a steep
one is nearly free.  We shall see that dependence explicitly in
Table~\ref{tab:11-omega}.

\paragraph{Single-particle states and magic numbers.}
The eigenstates of $\hat{H}_0$ for one particle in two dimensions are labelled
by a radial quantum number $n=0,1,2,\dots$ and an angular-momentum projection
$m=0,\pm1,\pm2,\dots$, with energies
\begin{equation}
  \varepsilon_{nm} = \hbar\omega\left(2n+|m|+1\right).
  \label{eq:11-spenergies}
\end{equation}
States with the same $N_s=2n+|m|$ are degenerate, so the spectrum falls into
shells of degeneracy $2(N_s+1)$ including spin.  Filling shells completely
gives the closed-shell electron numbers
\begin{equation}
  N = 2,\;6,\;12,\;20,\;30,\;42,\dots
  \label{eq:11-magic}
\end{equation}
the \emph{magic numbers} of the quantum dot, listed shell by shell in
Table~\ref{tab:11-shells}.  They play exactly the role that the closed shells
of atomic and nuclear physics do: adding or removing an electron at a magic
number costs more energy than elsewhere, and a closed-shell dot has a
non-degenerate ground state which a single Slater determinant can describe
reasonably.  All the machinery of Chapters~\ref{chap:hartrefock}
to~\ref{chap:cctheory} assumes such a reference, so we work at $N=2$ and
$N=6$.

\begin{table}[h]
\centering
\renewcommand{\arraystretch}{1.25}
\setlength{\tabcolsep}{10pt}
\begin{tabular}{clcc}
\toprule
$N_s$ & states $(n,m)$ & degeneracy & cumulative $N$\\
\midrule
$0$ & $(0,0)$ & $2$ & $\mathbf{2}$\\
$1$ & $(0,-1),\ (0,+1)$ & $4$ & $\mathbf{6}$\\
$2$ & $(0,-2),\ (1,0),\ (0,+2)$ & $6$ & $\mathbf{12}$\\
$3$ & $(0,\pm3),\ (1,\pm1)$ & $8$ & $\mathbf{20}$\\
$4$ & $(0,\pm4),\ (1,\pm2),\ (2,0)$ & $10$ & $\mathbf{30}$\\
$5$ & $(0,\pm5),\ (1,\pm3),\ (2,\pm1)$ & $12$ & $\mathbf{42}$\\
\bottomrule
\end{tabular}
\caption{Shell structure of the two-dimensional harmonic oscillator.  Shell
$N_s=2n+|m|$ has $N_s+1$ spatial states and twice that many spin-orbitals.
The cumulative column is the sequence of magic
numbers~(\ref{eq:11-magic}).  The calculations of this chapter use all six
shells, that is $42$ spin-orbitals.}
\label{tab:11-shells}
\end{table}

%----------------------------------------------------------------------
\section{The single-particle basis}
\label{sec:qdbasis}

\index{quantum dot!single-particle basis}
In Cartesian coordinates the two-dimensional oscillator separates into two
one-dimensional ones and the eigenfunctions are products of Hermite
polynomials.  That is correct but inconvenient: the rotational symmetry of the
trap is then invisible, and it is precisely that symmetry which will cut the
number of non-zero Coulomb matrix elements by two thirds.  In polar
coordinates $x=r\cos\theta$, $y=r\sin\theta$ the angular momentum is manifest.

Writing $\psi(r,\theta)=R(r)\,Y(\theta)$, the angular factor is
\begin{equation}
  Y(\theta) = \frac{1}{\sqrt{2\pi}}\,e^{\imath m\theta},
  \label{eq:11-angular}
\end{equation}
and single-valuedness, $\psi(r,\theta+2\pi)=\psi(r,\theta)$, forces $m$ to be
an integer.  The radial equation is solved by associated Laguerre polynomials,
\begin{equation}
  R_{nm}(r) = \sqrt{\frac{2\,n!}{(n+|m|)!}}\;
    \beta^{\frac12(|m|+1)}\,
    r^{|m|}\,e^{-\beta r^2/2}\,
    L_n^{|m|}\!\left(\beta r^2\right),
  \qquad \beta = \frac{m_e\omega}{\hbar},
  \label{eq:11-radial}
\end{equation}
so that the complete single-particle eigenfunction is
\begin{equation}
  \boxed{\;
  \psi_{nm}(r,\theta)
   = \sqrt{\frac{n!}{\pi\,(n+|m|)!}}\;
     \beta^{\frac12(|m|+1)}\,
     r^{|m|}\,e^{-\beta r^2/2}\,
     L_n^{|m|}\!\left(\beta r^2\right)e^{\imath m\theta} \;}
  \label{eq:11-eigenfunction}
\end{equation}
with the eigenvalue~(\ref{eq:11-spenergies}).  The Cartesian and polar
labellings agree, $n_x+n_y=2n+|m|$, as they must.

This basis is not the only possible one, and it is not obviously the best.  It
is the eigenbasis of the non-interacting problem, which makes $\hat{H}_0$
diagonal and gives a natural ordering by shell, but it is built from Gaussians
and therefore cannot reproduce the cusp that the Coulomb interaction imposes
on the exact wave function at $r_{12}=0$.  That single fact governs how the
results below converge, and it is worth keeping in mind when reading
Table~\ref{tab:11-twoelectron}.

%----------------------------------------------------------------------
\section{The Coulomb matrix element}
\label{sec:qdcoulomb}

\index{Coulomb matrix element}
Everything now hinges on
\begin{equation}
  \bracket{pq}{\hat{v}}{rs}
   = \left\langle (n_pm_p)(n_qm_q)\left\vert
       \frac{1}{|\bm{r}_1-\bm{r}_2|}
     \right\vert (n_rm_r)(n_sm_s)\right\rangle ,
  \label{eq:11-coulombelement}
\end{equation}
a four-dimensional integral over two planes.  It can be done in closed form.
The derivation is due to Anisimova and Matulis and consists of expanding the
Laguerre polynomials binomially, doing the angular integrals with the
selection rule below, and recognising the radial integrals as Gamma functions;
the result is a finite sum of products of factorials, which
\texttt{quantumdot.py} evaluates in logarithms so that nothing overflows.

Two exact statements about Eq.~(\ref{eq:11-coulombelement}) do more for us
than the closed form itself.

\paragraph{Angular momentum is conserved.}
The Coulomb interaction depends only on $|\bm{r}_1-\bm{r}_2|$ and is therefore
invariant under a simultaneous rotation of both particles.  Hence
\begin{equation}
  \boxed{\;m_p+m_q = M = m_r+m_s\;}
  \label{eq:11-mconservation}
\end{equation}
and the matrix element vanishes identically otherwise.  This is the
single most valuable fact in the implementation: of the $M_s^4$ index
combinations for $M_s$ spatial states, the great majority are eliminated
before any work is done.  For the six-shell basis, $21^4=194\,481$
combinations reduce to $14\,703$ non-zero ones.

\paragraph{The frequency dependence is a scale factor.}
Lengths scale as $\omega^{-1/2}$ and the interaction as one over a length, so
\begin{equation}
  \bracket{pq}{\hat{v}}{rs}_{\omega}
   = \sqrt{\hbar\omega}\;\bracket{pq}{\hat{v}}{rs}_{\omega=1}.
  \label{eq:11-omegascaling}
\end{equation}
The expensive integrals need therefore be computed only once, and the whole
$\omega$ dependence of Table~\ref{tab:11-omega} costs nothing extra.  It also
explains the physics: since $\hat{H}_0$ scales as $\omega$ and $\hat{V}$ as
$\sqrt{\omega}$, the ratio of interaction to kinetic energy goes as
$\omega^{-1/2}$, and correlation is a small-$\omega$ phenomenon.

\paragraph{Spin.}
The closed form is a spatial integral; the spin has to be put in by hand.  The
antisymmetrised matrix element that all our machinery expects is
\begin{equation}
  \bracket{pq}{\hat{v}}{rs}_{\mathrm{AS}}
   = \bracket{pq}{\hat{v}}{rs}\,
     \delta_{\sigma_p\sigma_r}\delta_{\sigma_q\sigma_s}
   - \bracket{pq}{\hat{v}}{sr}\,
     \delta_{\sigma_p\sigma_s}\delta_{\sigma_q\sigma_r},
  \label{eq:11-antisym}
\end{equation}
the direct term requiring each particle to keep its spin and the exchange term
requiring the two to swap.  Both together conserve $S_z$.  It is easy to get
the second delta pair wrong -- note that the roles of $r$ and $s$ are
interchanged in the exchange term -- and the resulting error is subtle, since
the calculation still runs and still converges.  The checks below exist for
exactly that reason.

\paragraph{Verification.}
Three tests catch almost every implementation error.  The lowest element is
known analytically,
\begin{equation}
  \bracket{00;00}{\hat{v}}{00;00}
   = \sqrt{\frac{\pi}{2}}\,\sqrt{\hbar\omega}
   = 1.253314137316\ \sqrt{\hbar\omega},
  \label{eq:11-lowestelement}
\end{equation}
which \texttt{quantumdot.py} reproduces to twelve digits.  The scaling
(\ref{eq:11-omegascaling}) can be checked by computing the same element at
several frequencies and dividing.  And the element must satisfy
\begin{equation}
  \bracket{pq}{\hat{v}}{rs} = \bracket{qp}{\hat{v}}{sr}
   = \bracket{rs}{\hat{v}}{pq},
  \label{eq:11-symmetries}
\end{equation}
particle exchange and hermiticity, both of which hold to machine precision.
Equation~(\ref{eq:11-lowestelement}) is also the first-order perturbative
estimate of the two-electron ground-state energy,
$E\approx2\hbar\omega+\sqrt{\pi/2}\sqrt{\hbar\omega}$, which at
$\hbar\omega=1$ gives $3.253314$ -- the minimal-basis Hartree-Fock energy of
the next section.

%----------------------------------------------------------------------
\section{Hartree-Fock}
\label{sec:qdhf}

\index{Hartree-Fock theory!for quantum dots}
The Hartree-Fock equations of Section~\ref{sec:hfcoefficients} are used here
without change.  We expand the Hartree-Fock orbitals in the oscillator basis,
\begin{equation}
  \psi_{\alpha} = \sum_p C_{p\alpha}\,\phi_p ,
  \label{eq:11-expansion}
\end{equation}
and the variational condition becomes the Roothaan eigenvalue problem
$\bm{F}\bm{C}=\bm{C}\bm{\varepsilon}$ with
\begin{equation}
  F_{pq} = \bracket{p}{\hat{h}_0}{q}
    + \sum_{rs}\rho_{rs}\,\bracket{pr}{\hat{v}}{qs}_{\mathrm{AS}},
  \qquad
  \rho_{rs} = \sum_{a=1}^{N}C_{ra}C_{sa} ,
  \label{eq:11-fock}
\end{equation}
solved by the self-consistent iteration of Section~\ref{sec:scf}.  In the
oscillator basis $\bracket{p}{\hat{h}_0}{q}=\varepsilon_p^{(0)}\delta_{pq}$
with $\varepsilon^{(0)}_p$ from Eq.~(\ref{eq:11-spenergies}), so the only
non-trivial ingredient is the two-body matrix built in the previous section.
The total energy is
\begin{equation}
  E_{\mathrm{HF}}
   = \sum_{a=1}^{N}\varepsilon^{(0)}_a
   + \frac12\sum_{a,b=1}^{N}\bracket{ab}{\hat{v}}{ab}_{\mathrm{AS}} ,
  \label{eq:11-hfenergy}
\end{equation}
the factor $\tfrac12$ removing the double counting of the interaction.

\paragraph{Convergence with the basis.}
Table~\ref{tab:11-hfconvergence} follows the energy as shells are added.  The
sequence is monotone, as it must be: enlarging the basis gives the
Hartree-Fock state more variational freedom and can only lower the energy.
Two features are worth noting.  The energy is unchanged when certain shells
are added -- from one shell to two for $N=2$, and again from three to four --
for a reason worked out in Exercise~\ref{ex:11-hfcheck}.  And convergence is
slow at the end: going from $30$ to $42$ orbitals still moves the six-electron
energy by $0.028$.  That is the Gaussian basis struggling with the Coulomb
cusp, and it is a much larger effect than anything the many-body method will
contribute.

\begin{table}[h]
\centering
\renewcommand{\arraystretch}{1.25}
\setlength{\tabcolsep}{9pt}
\begin{tabular}{rrllll}
\toprule
shells & orbitals & $N=2$: $E_{\mathrm{HF}}$ & $\Delta E$ &
$N=6$: $E_{\mathrm{HF}}$ & $\Delta E$\\
\midrule
$0$ & $2$  & $3.25331414$ & --- & --- & ---\\
$1$ & $6$  & $3.25331414$ & $+0.000000$ & $22.21981284$ & ---\\
$2$ & $12$ & $3.16269135$ & $-0.090623$ & $21.59319848$ & $-0.626614$\\
$3$ & $20$ & $3.16269135$ & $+0.000000$ & $20.76691943$ & $-0.826279$\\
$4$ & $30$ & $3.16192140$ & $-0.000770$ & $20.74840225$ & $-0.018517$\\
$5$ & $42$ & $3.16192140$ & $+0.000000$ & $20.72025707$ & $-0.028145$\\
\bottomrule
\end{tabular}
\caption{Hartree-Fock energy as a function of the number of oscillator shells,
$\hbar\omega=1$.  $\Delta E$ is the change from the previous row.  Computed by
\texttt{quantumdot.py}.}
\label{tab:11-hfconvergence}
\end{table}

\paragraph{The single-particle spectrum.}
The converged Hartree-Fock energies for six electrons in the full basis are
\begin{equation}
  \varepsilon_{1,2} = 4.54225444,
  \qquad
  \varepsilon_{3,4,5,6} = 5.30056289,
  \label{eq:11-hfspectrum}
\end{equation}
against the bare oscillator values $1$ and $2$.  The mean field has pushed
every level up by more than three units -- the Coulomb repulsion of the other
electrons -- and the degeneracies of Table~\ref{tab:11-shells} survive
untouched, which is a useful check that the self-consistent solution has not
broken the rotational symmetry.  By Koopmans' theorem the ionisation energy of
the dot is $-\varepsilon_6=-5.30056289$; for two electrons the corresponding
number is $-\varepsilon_2=-2.12246475$.

%----------------------------------------------------------------------
\section{Many-body perturbation theory}
\label{sec:qdmbpt}

\index{many-body perturbation theory!for quantum dots}
With a Hartree-Fock reference in hand, the natural partition is the
M{\o}ller-Plesset one of Chapter~\ref{chap:mbpt}: the unperturbed operator is
the sum of Fock operators and the perturbation is the \emph{fluctuation
potential}, the difference between the true interaction and the mean field it
generates,
\begin{equation}
  \hat{H}_0^{\mathrm{MP}} = \sum_{i=1}^{N}\hat{f}(i),
  \qquad
  \hat{V}' = \hat{V} - \sum_{i=1}^{N}\hat{u}^{\mathrm{HF}}(i).
  \label{eq:11-mppartition}
\end{equation}
The first-order energy then vanishes by construction,
$E^{(0)}+E^{(1)}=E_{\mathrm{HF}}$, so the leading correction is second order.
Brillouin's theorem removes the singly excited determinants,
$\bracket{\Phi_a^r}{\hat{V}'}{\Phi_0}=0$, and only doubles survive:
\begin{equation}
  \boxed{\;
  E_{\mathrm{MP2}}
   = \frac14\sum_{ab}^{\mathrm{occ}}\sum_{rs}^{\mathrm{virt}}
     \frac{\left\vert\bracket{ab}{\hat{v}}{rs}_{\mathrm{AS}}\right\vert^2}
          {\varepsilon_a+\varepsilon_b-\varepsilon_r-\varepsilon_s} \;}
  \label{eq:11-mp2}
\end{equation}
The denominator is negative, so MP2 always lowers the energy.  The matrix
elements are needed in the Hartree-Fock basis, obtained from the oscillator
ones by the four-index transformation
\begin{equation}
  \bracket{\alpha\beta}{\hat{v}}{\gamma\delta}
   = \sum_{pqrs}C_{p\alpha}C_{q\beta}\,
     \bracket{pq}{\hat{v}}{rs}\,C_{r\gamma}C_{s\delta},
  \label{eq:11-transformation}
\end{equation}
carried out in two half-steps so that the cost is $\bigO(M^5)$ rather than
$\bigO(M^8)$.

For the full basis at $\hbar\omega=1$ the result is
\begin{equation}
  E_{\mathrm{MP2}}(N=2) = -0.13488329,
  \qquad
  E_{\mathrm{MP2}}(N=6) = -0.41769581,
  \label{eq:11-mp2results}
\end{equation}
that is $11.6\%$ and $3.9\%$ of the Hartree-Fock interaction energy
$E_{\mathrm{HF}}-E_0$.  The smaller fraction for six electrons is not a
statement that the six-electron dot is less correlated; it is that the mean
field is doing more of the work, since each electron feels five others rather
than one.

%----------------------------------------------------------------------
\section{Coupled cluster}
\label{sec:qdcc}

\index{coupled cluster!for quantum dots}
The coupled-cluster equations of Chapter~\ref{chap:cctheory} are applied here
verbatim: the same amplitude equations, the same
Stanton-Gauss-Watts-Bartlett intermediates, the same solver.  Only the
matrices change.  That is the whole point of writing the equations in a
spin-orbital basis with an antisymmetrised interaction -- the method never
learns what the interaction is.

\paragraph{A warning about the singles.}
It is often said that for a closed-shell system in a canonical Hartree-Fock
basis the singles amplitudes vanish by Brillouin's theorem, so that CCSD
reduces to CCD.  This is not true, and it is worth being explicit about why.
Brillouin's theorem states that $f_{ai}=0$, and that is what kills the
first-order singles contribution in perturbation theory.  But the CCSD singles
equation is not homogeneous in $t_1$: it contains the terms
\begin{equation}
  \tfrac12\sum_{mef}t_{im}^{ef}\,\bracket{ma}{\hat{v}}{ef}_{\mathrm{AS}}
  \qquad\text{and}\qquad
  -\tfrac12\sum_{men}t_{mn}^{ae}\,\bracket{nm}{\hat{v}}{ei}_{\mathrm{AS}},
  \label{eq:11-singlesdriver}
\end{equation}
which are built from the \emph{doubles} and survive at $t_1=0$.  The singles
are therefore driven by the doubles, they are small but not zero, and they
describe the relaxation of the orbitals in response to the correlation the
doubles have introduced.  For the quantum dots computed here
$\max|t_i^a| = 5.7\times10^{-3}$ at $N=2$ and $1.0\times10^{-2}$ at $N=6$, and
CCSD lies correspondingly a little below CCD.  A calculation that reports
$t_1=0$ identically has a bug in its singles equation, and
Exercise~\ref{ex:11-singles} gives a cheap test that finds it.

\paragraph{Results.}
Table~\ref{tab:11-hierarchy} collects the whole hierarchy in the full
$42$-orbital basis.  The ordering
\begin{equation}
  E_{\mathrm{HF}} > E_{\mathrm{HF}}+E_{\mathrm{MP2}}
   > E_{\mathrm{HF}}+E_{\mathrm{CCD}} > E_{\mathrm{HF}}+E_{\mathrm{CCSD}}
  \label{eq:11-ordering}
\end{equation}
holds for both electron numbers, and each step has a clear meaning: MP2 is the
leading pair correlation, CCD resums the ladders and rings to all orders in
$\hat{T}_2$, and CCSD adds the orbital relaxation just discussed.

\begin{table}[h]
\centering
\renewcommand{\arraystretch}{1.25}
\setlength{\tabcolsep}{14pt}
\begin{tabular}{lll}
\toprule
quantity & $N=2$ & $N=6$\\
\midrule
non-interacting $E_0$ & $\phantom{-}2.00000000$ & $\phantom{-}10.00000000$\\
Hartree-Fock $E_{\mathrm{HF}}$ & $\phantom{-}3.16192140$ &
  $\phantom{-}20.72025707$\\
\quad interaction $E_{\mathrm{HF}}-E_0$ & $\phantom{-}1.16192140$ &
  $\phantom{-}10.72025707$\\
\midrule
$E_{\mathrm{MP2}}$ & $-0.13488329$ & $-0.41769581$\\
\quad total & $\phantom{-}3.02703812$ & $\phantom{-}20.30256126$\\
$E_{\mathrm{CCD}}$ & $-0.14799908$ & $-0.44624450$\\
\quad total & $\phantom{-}3.01392232$ & $\phantom{-}20.27401257$\\
$E_{\mathrm{CCSD}}$ & $-0.14829527$ & $-0.44701044$\\
\quad total & $\phantom{-}3.01362613$ & $\phantom{-}20.27324663$\\
\bottomrule
\end{tabular}
\caption{The hierarchy of methods for the two- and six-electron quantum dot at
$\hbar\omega=1$, in the basis of six oscillator shells ($42$ spin-orbitals).
The self-consistent field converges in $15$ and $24$ iterations, the
coupled-cluster equations in $23$ and $28$.  Computed by
\texttt{quantumdot.py}.}
\label{tab:11-hierarchy}
\end{table}

%----------------------------------------------------------------------
\section{Comparison with the exact answer}
\label{sec:qdexact}

Tables of energies produced by approximate methods are worth little without
something to compare them with.  For this system we have two benchmarks.

\paragraph{Two electrons: exact within the basis, and exact in the limit.}
With $N=2$ the determinant space has dimension $\binom{M}{2}$, which is $861$
for $42$ spin-orbitals and can be diagonalised directly.  Moreover singles and
doubles \emph{exhaust} the excitations available to two particles, so CCSD
must be exact in any basis.  Table~\ref{tab:11-twoelectron} confirms both
statements: CCSD and full diagonalisation agree to twelve digits at every
basis size.  This is a stronger test of the coupled-cluster solver of
Chapter~\ref{chap:cctheory} than anything the pairing model could provide,
because here the interaction has no special structure to hide behind.

The second benchmark is analytic.  Taut showed that the relative motion of two
electrons in a parabolic trap is exactly solvable for a discrete set of
frequencies, and at $\hbar\omega=1$ in our units the ground state is exactly
$E=3$.  The last column of Table~\ref{tab:11-twoelectron} approaches that
value from above, but slowly: at $42$ orbitals the remaining error is
$0.0136$, comparable with the entire difference between MP2 and CCSD.  The
many-body problem has been solved exactly; the basis-set problem has not.  For
this system the truncation of the single-particle space, not the truncation of
the cluster operator, is the dominant error -- which is the opposite of the
situation in the pairing model, and a lesson worth taking from a realistic
Hamiltonian.

\begin{table}[h]
\centering
\renewcommand{\arraystretch}{1.25}
\setlength{\tabcolsep}{6pt}
\small
\begin{tabular}{rrrlllll}
\toprule
shells & orb. & $\dim$ & $E_{\mathrm{HF}}$ &
$+E_{\mathrm{MP2}}$ & $+E_{\mathrm{CCD}}$ & $+E_{\mathrm{CCSD}}$ &
$E_{\mathrm{FCI}}$\\
\midrule
$1$ & $6$  & $15$  & $3.25331414$ & $3.17856150$ & $3.15232801$ &
  $3.15232801$ & $3.15232801$\\
$2$ & $12$ & $66$  & $3.16269135$ & $3.05797643$ & $3.03904782$ &
  $3.03860458$ & $3.03860458$\\
$3$ & $20$ & $190$ & $3.16269135$ & $3.04440370$ & $3.02527309$ &
  $3.02523058$ & $3.02523058$\\
$4$ & $30$ & $435$ & $3.16192140$ & $3.03341846$ & $3.01794371$ &
  $3.01760623$ & $3.01760623$\\
$5$ & $42$ & $861$ & $3.16192140$ & $3.02703812$ & $3.01392232$ &
  $3.01362613$ & $3.01362613$\\
\midrule
$\infty$ & --- & --- & --- & --- & --- & --- & $3.00000000$\\
\bottomrule
\end{tabular}
\caption{Two electrons at $\hbar\omega=1$.  CCSD agrees with full
diagonalisation to $3\times10^{-13}$ in every basis, as it must for two
particles.  The last row is Taut's analytic result for the complete basis.
Computed by \texttt{quantumdot.py}.}
\label{tab:11-twoelectron}
\end{table}

\paragraph{Six electrons: exact in a small basis.}
For six electrons the full space is out of reach at $42$ orbitals --
$\binom{42}{6}=5\,245\,786$ determinants -- but in the two-shell basis it is
only $\binom{12}{6}=924$, and there the truncation of the cluster operator
genuinely matters.  Table~\ref{tab:11-sixelectron} is therefore the one
comparison in this book where every method is measured against the exact
answer on a realistic interaction in a system that is not trivially solvable.
The last column is the fraction of the exact correlation energy recovered.

\begin{table}[h]
\centering
\renewcommand{\arraystretch}{1.25}
\setlength{\tabcolsep}{11pt}
\begin{tabular}{llrr}
\toprule
method & energy & error & \% of $E_{\mathrm{corr}}$\\
\midrule
Hartree-Fock          & $21.59319848$ & $1.7\times10^{-1}$ & $\phantom{00}0.00$\\
MP2                   & $21.43344056$ & $1.3\times10^{-2}$ & $\phantom{0}92.55$\\
CCD                   & $21.42380972$ & $3.2\times10^{-3}$ & $\phantom{0}98.13$\\
CCSD                  & $21.42323158$ & $2.6\times10^{-3}$ & $\phantom{0}98.47$\\
UCCSD                 & $21.42196510$ & $1.4\times10^{-3}$ & $\phantom{0}99.20$\\
\midrule
exact diagonalisation & $21.42058830$ & --- & $100.00$\\
\bottomrule
\end{tabular}
\caption{Six electrons at $\hbar\omega=1$ in the two-shell basis, $12$
spin-orbitals, determinant space of dimension $924$.  The correlation energy
is $E_{\mathrm{corr}}=-0.17261018$.  The unitary result is evaluated at the
converged CCSD amplitudes rather than at its own variational optimum, and is
still the best of the four.  Computed by \texttt{quantumdot.py}.}
\label{tab:11-sixelectron}
\end{table}

\paragraph{Dependence on the confinement.}
Table~\ref{tab:11-omega} varies the trap frequency.  The last column, the
correlation energy as a fraction of the Hartree-Fock interaction energy, falls
by a factor of three between $\hbar\omega=0.25$ and $\hbar\omega=4$, which is
the $\omega^{-1/2}$ scaling of Eq.~(\ref{eq:11-omegascaling}) doing its work.
A shallow trap is a strongly correlated system: the electrons have room to
avoid one another, the mean field describes them badly, and the correlation
methods have more to do.  In the opposite limit the dot becomes a set of
nearly free particles in a very steep well and Hartree-Fock is almost exact.
This is the same story the pairing model told through $g$ in
Table~\ref{tab:pairing}, told now by a physical parameter that an
experimentalist can actually tune.

\begin{table}[h]
\centering
\renewcommand{\arraystretch}{1.25}
\setlength{\tabcolsep}{8pt}
\small
\begin{tabular}{rrllll}
\toprule
$\hbar\omega$ & $E_0$ & $E_{\mathrm{HF}}$ & $E_{\mathrm{MP2}}$ &
$E_{\mathrm{CCSD}}$ & ratio\\
\midrule
\multicolumn{6}{l}{\emph{$N=2$}}\\
$0.25$ & $0.5$  & $\phantom{0}1.04628424$ & $-0.10380945$ & $-0.11288484$ &
  $-0.2066$\\
$0.50$ & $1.0$  & $\phantom{0}1.79974794$ & $-0.12036392$ & $-0.13249076$ &
  $-0.1657$\\
$1.00$ & $2.0$  & $\phantom{0}3.16192140$ & $-0.13488329$ & $-0.14829527$ &
  $-0.1276$\\
$2.00$ & $4.0$  & $\phantom{0}5.67728232$ & $-0.14693248$ & $-0.15989645$ &
  $-0.0953$\\
$4.00$ & $8.0$  & $10.40864632$ & $-0.15651039$ & $-0.16789538$ & $-0.0697$\\
\midrule
\multicolumn{6}{l}{\emph{$N=6$}}\\
$0.25$ & $2.5$  & $\phantom{0}7.39178633$ & $-0.32952382$ & $-0.35507942$ &
  $-0.0726$\\
$0.50$ & $5.0$  & $12.27149922$ & $-0.37861963$ & $-0.40804475$ & $-0.0561$\\
$1.00$ & $10.0$ & $20.72025707$ & $-0.41769581$ & $-0.44701044$ & $-0.0417$\\
$2.00$ & $20.0$ & $35.67233283$ & $-0.44700367$ & $-0.47274020$ & $-0.0302$\\
$4.00$ & $40.0$ & $62.74116274$ & $-0.46768316$ & $-0.48840005$ & $-0.0215$\\
\bottomrule
\end{tabular}
\caption{Dependence on the trap frequency, in the full $42$-orbital basis.
The last column is $E_{\mathrm{CCSD}}/(E_{\mathrm{HF}}-E_0)$.  For $N=2$ the
CCSD column is exact within the basis, so
$E_{\mathrm{HF}}+E_{\mathrm{CCSD}}$ is the full configuration-interaction
energy at every frequency.  Computed by \texttt{quantumdot.py}.}
\label{tab:11-omega}
\end{table}

%----------------------------------------------------------------------
\section{Unitary coupled cluster}
\label{sec:qducc}

\index{coupled cluster!unitary!for quantum dots}
The unitary ansatz of Section~\ref{sec:ucc} carries over unchanged.  With
\begin{equation}
  \hat{\sigma} = \hat{T}-\hat{T}^{\dagger},
  \qquad
  \bket{\Psi(\bm{t})} = e^{\hat{\sigma}(\bm{t})}\bket{\Phi_0},
  \qquad
  E(\bm{t}) = \bracket{\Psi(\bm{t})}{\hat{H}}{\Psi(\bm{t})},
  \label{eq:11-ucc}
\end{equation}
the energy is a genuine expectation value in a normalised state and is
therefore variational: it lies above the exact energy for any amplitudes,
which the projected CCSD energy does not.  The price is that the
Baker-Campbell-Hausdorff expansion no longer terminates, so there are no
closed amplitude equations and the energy has to be evaluated by
exponentiating $\hat{\sigma}$.

\paragraph{Two electrons: everything collapses.}
Table~\ref{tab:11-ucctwo} runs the unitary calculation for two electrons.  Two
things happen, both instructive.  UCCSD, CCSD and full diagonalisation
coincide, which is expected: with two particles singles and doubles span the
whole space and any ansatz that is exact there must give the exact energy.
Less obviously, UCCD coincides with CCD.  The reason is that for two particles
$\hat{T}_2\hat{T}_2$ annihilates the reference and so does
$\hat{T}_2^{\dagger}$, so $e^{\hat{T}_2-\hat{T}_2^{\dagger}}\bket{\Phi_0}$ and
$(1+\hat{T}_2)\bket{\Phi_0}$ sweep out the same two-dimensional manifold; the
variational minimum over it is the lowest eigenvalue in the space spanned by
the reference and the doubles, which is what the CCD equations return.  The
unitary and the standard ansatz differ only when the truncation genuinely
bites.

\begin{table}[h]
\centering
\renewcommand{\arraystretch}{1.25}
\setlength{\tabcolsep}{6pt}
\small
\begin{tabular}{rrrlllll}
\toprule
shells & orb. & ampl. & CCD & UCCD & CCSD & UCCSD & FCI\\
\midrule
$1$ & $6$  & $14$  & $3.15232801$ & $3.15232801$ & $3.15232801$ &
  $3.15232801$ & $3.15232801$\\
$2$ & $12$ & $65$  & $3.03904782$ & $3.03904782$ & $3.03860458$ &
  $3.03860458$ & $3.03860458$\\
$3$ & $20$ & $189$ & $3.02527309$ & $3.02527309$ & $3.02523058$ &
  $3.02523058$ & $3.02523058$\\
\bottomrule
\end{tabular}
\caption{Unitary and standard coupled cluster for two electrons at
$\hbar\omega=1$.  The unitary energies are obtained by minimising
Eq.~(\ref{eq:11-ucc}) over all amplitudes, the standard ones by solving the
projected equations.  Computed by \texttt{quantumdot.py}.}
\label{tab:11-ucctwo}
\end{table}

\paragraph{Six electrons: the unitary ansatz wins.}
The six-electron row of Table~\ref{tab:11-sixelectron} is where the difference
appears.  There the truncation to singles and doubles is a real
approximation -- the exact state contains triples and quadruples -- and the
unitary ansatz recovers $99.2\%$ of the correlation energy against CCSD's
$98.5\%$, while remaining above the exact energy as the variational principle
requires.  What makes this worth emphasising is that the unitary number was
obtained \emph{without optimising}: the converged CCSD amplitudes were simply
substituted into Eq.~(\ref{eq:11-ucc}).  Even at borrowed amplitudes the
unitary ansatz is the better wave function, because $e^{\hat{\sigma}}$
generates the de-excitations that $e^{\hat{T}}$ cannot.

The cost of doing better is what pushes one towards a quantum computer.  A
proper variational optimisation of the six-electron case means minimising over
$261$ amplitudes, each energy evaluation requiring the exponential of a
$924\times924$ matrix, and both numbers grow exponentially with the size of
the system.  On a quantum computer the state $e^{\hat{\sigma}}\bket{\Phi_0}$
is prepared by a circuit whose depth grows polynomially, and the energy is
measured rather than computed.

\paragraph{Trotterisation.}
A circuit cannot apply $e^{\hat{\sigma}}$ in one step; it applies the
generators one at a time, and the resulting splitting error is the Trotter
error of Section~\ref{sec:trotter}.  Table~\ref{tab:11-trotter} measures it
for the two-electron dot at the CCSD amplitudes.  The error halves when the
number of Trotter steps doubles -- the $\bigO(1/n)$ rate of
Eq.~(\ref{eq:6-trotter}) -- and $n$ is the circuit depth.  On hardware the
Hamiltonian would first be mapped to Pauli strings by the Jordan-Wigner
transformation of Section~\ref{sec:jordanwigner}, each string contributing one
rotation to the circuit; here the matrices are simply exponentiated.

\begin{table}[h]
\centering
\renewcommand{\arraystretch}{1.25}
\setlength{\tabcolsep}{14pt}
\begin{tabular}{rlll}
\toprule
Trotter steps $n$ & energy & error & ratio\\
\midrule
$1$  & $3.03863183$ & $1.04\times10^{-5}$ & ---\\
$2$  & $3.03863789$ & $4.38\times10^{-6}$ & $2.38$\\
$4$  & $3.03864044$ & $1.83\times10^{-6}$ & $2.40$\\
$8$  & $3.03864146$ & $8.10\times10^{-7}$ & $2.26$\\
$16$ & $3.03864189$ & $3.77\times10^{-7}$ & $2.15$\\
\midrule
$\infty$ & $3.03864227$ & --- & ---\\
\bottomrule
\end{tabular}
\caption{First-order Trotterisation of the unitary ansatz, two electrons in
$12$ orbitals at the converged CCSD amplitudes.  The error is measured against
the untrotterised exponential and falls off as $1/n$, the ratio approaching
$2$ from above.  Computed by \texttt{quantumdot.py}.}
\label{tab:11-trotter}
\end{table}

\begin{notebox}
\textbf{Why this system matters.}  The quantum dot is the smallest realistic
many-fermion system on which the whole apparatus of this book can be run and
checked.  It is a genuine Coulomb problem, so the matrix elements and the
basis-set convergence are those of quantum chemistry and nuclear physics; it
is small enough that exact diagonalisation is available for two electrons in
any basis and for six in a small one; and it has an analytic answer at special
frequencies.  It is also the system on which the stochastic methods of the
later chapters will be tested, so the numbers computed here are the reference
against which variational and diffusion Monte Carlo will be measured.
\end{notebox}

%----------------------------------------------------------------------
\section{What the realistic system taught us}
\label{sec:qdsummary}

The methods survived the transition from a one-parameter model to a real
interaction without modification.  The same Hartree-Fock routine, the same
MP2 sum, the same coupled-cluster solver and the same unitary ansatz that were
developed and validated on the pairing model of Chapter~\ref{chap:systems}
produce Tables~\ref{tab:11-hierarchy} to~\ref{tab:11-trotter} when handed a
different pair of matrices.  That portability is not an accident; it is what
the spin-orbital formulation of Chapter~\ref{chap:secondquant} and the
antisymmetrised matrix element were for.

Three things did change, and they are the reasons for doing a realistic
calculation at all.

\begin{itemize}
\item \textbf{The basis is now an approximation.}  In the pairing model the
  single-particle space was finite and the only error was the many-body
  method.  Here the oscillator basis is truncated, and for two electrons the
  truncation error at $42$ orbitals is $0.0136$ against a CCSD-minus-MP2
  difference of $0.0134$.  The two errors are comparable, and reporting one
  without the other is meaningless.
\item \textbf{The hierarchy is quantitative, not qualitative.}
  Table~\ref{tab:11-sixelectron} puts numbers on what each method buys:
  MP2 recovers $92.6\%$ of the correlation energy, CCD $98.1\%$, CCSD
  $98.5\%$ and the unitary ansatz $99.2\%$.  The steps get smaller and the
  costs get larger, which is the shape of the whole subject.
\item \textbf{Folklore does not survive contact with a computation.}  The
  singles amplitudes do not vanish for a closed-shell Hartree-Fock reference,
  CCSD is not equal to CCD, and a calculation reporting otherwise is reporting
  a bug.  The check that catches it is cheap: for two electrons CCSD must
  equal full diagonalisation, and it does, to twelve digits.
\end{itemize}

The Hamiltonian~(\ref{eq:11-hamiltonian}) returns later in the book.  In the
stochastic chapters it is treated by variational and diffusion Monte Carlo,
which sample the many-body wave function directly rather than expanding it in
determinants, and which therefore have no basis-set error at all -- their
error is statistical, and of a completely different nature.  Comparing the two
families on the same system, with the numbers of this chapter as the
reference, is the cleanest way to see what each is good for.

%=============================================================================
\section{Exercises}
%=============================================================================

\subsection*{Exercise 1: the shell structure}
\label{ex:11-shells}
\addcontentsline{toc}{subsection}{Exercise 1: the shell structure}

\begin{enumerate}
  \item[a)] Show that the number of spatial states with $2n+|m|=N_s$ is
  $N_s+1$, and hence that the magic numbers are $N=(N_s+1)(N_s+2)$.
  \item[b)] Verify that the Cartesian and polar labellings give the same
  spectrum.
  \item[c)] Why is $N=4$ not a magic number, and what would go wrong if we
  attempted a Hartree-Fock calculation there?
\end{enumerate}

\paragraph{Answer to a).}
For fixed $N_s$ the allowed radial quantum numbers are
$n=0,1,\dots,\lfloor N_s/2\rfloor$ and $|m|=N_s-2n$.  Each $n$ with $|m|>0$
contributes two states, $m=\pm|m|$, and the single $n=N_s/2$ with $m=0$
(present only for even $N_s$) contributes one.  Counting gives $N_s+1$ spatial
states, or $2(N_s+1)$ with spin.  The cumulative sum is
$\sum_{k=0}^{N_s}2(k+1)=(N_s+1)(N_s+2)$, giving $2,6,12,20,30,42$.

\paragraph{Answer to b).}
The Cartesian energy is $\hbar\omega(n_x+n_y+1)$ and the polar one is
$\hbar\omega(2n+|m|+1)$, so the two agree whenever $n_x+n_y=2n+|m|$.  For a
given shell the Cartesian count is the number of pairs $(n_x,n_y)$ with
$n_x+n_y=N_s$, which is $N_s+1$ -- the same as in a).  The two bases are
related by a unitary transformation within each shell, and either may be used;
the polar one is chosen because it diagonalises $\hat{L}_z$.

\paragraph{Answer to c).}
$N=4$ falls inside the $N_s=1$ shell, which holds four spin-orbitals of which
only two would be occupied.  The reference determinant is then not unique --
any choice of two of the four degenerate orbitals gives the same unperturbed
energy -- and the ground state of the interacting system is degenerate or
nearly so.  A single-determinant reference is a poor starting point, the
self-consistent iteration may fail to converge, and the perturbation
denominators $\varepsilon_a+\varepsilon_b-\varepsilon_r-\varepsilon_s$ can
vanish.  This is the open-shell problem, and it needs multireference methods.

\subsection*{Exercise 2: the Coulomb matrix element}
\label{ex:11-coulomb}
\addcontentsline{toc}{subsection}{Exercise 2: the Coulomb matrix element}

\begin{enumerate}
  \item[a)] Prove the selection rule~(\ref{eq:11-mconservation}) from the
  angular integrals alone.
  \item[b)] Derive the scaling~(\ref{eq:11-omegascaling}).
  \item[c)] Estimate what fraction of the $M_s^4$ index combinations survives
  the selection rule for a basis of $M_s$ spatial states, and compare with the
  measured $14\,703$ out of $194\,481$ for six shells.
\end{enumerate}

\paragraph{Answer to a).}
Write $\bm{r}_1$ and $\bm{r}_2$ in polar coordinates.  The interaction
$1/|\bm{r}_1-\bm{r}_2|$ depends on $\theta_1$ and $\theta_2$ only through
$\theta_1-\theta_2$, so a Fourier expansion in that difference gives terms
$e^{\imath k(\theta_1-\theta_2)}$.  The angular part of the matrix element is
then
\[
  \int_0^{2\pi}\!\!\!\int_0^{2\pi}
    e^{-\imath m_p\theta_1}e^{-\imath m_q\theta_2}\,
    e^{\imath k(\theta_1-\theta_2)}\,
    e^{\imath m_r\theta_1}e^{\imath m_s\theta_2}
    \,d\theta_1\,d\theta_2 ,
\]
which is non-zero only when $m_r-m_p+k=0$ and $m_s-m_q-k=0$ simultaneously.
Adding the two conditions eliminates $k$ and leaves $m_p+m_q=m_r+m_s$.

\paragraph{Answer to b).}
The oscillator eigenfunctions~(\ref{eq:11-eigenfunction}) depend on $r$ only
through $\beta r^2$ with $\beta=m_e\omega/\hbar$, so the natural length in the
problem is $\beta^{-1/2}\propto\omega^{-1/2}$.  The matrix element of
$1/r_{12}$ is an inverse length times a dimensionless number that depends only
on the quantum numbers, and therefore scales as $\omega^{1/2}$.  Numerically,
dividing the element by $\sqrt{\hbar\omega}$ gives $1.253314137316$ at every
frequency, as \texttt{quantumdot.py} confirms.

\paragraph{Answer to c).}
Fix $m_p$, $m_q$ and $m_r$; then $m_s$ is determined, and the combination is
allowed only if that value lies in the basis.  Roughly one index in four is
fixed by the other three, so of order $M_s^3$ of the $M_s^4$ combinations
survive, a fraction $\bigO(1/M_s)$.  For $M_s=21$ this suggests about
$5\%$; the measured fraction is $14\,703/194\,481=7.6\%$, larger because the
determined $m_s$ often falls inside the basis rather than outside it,
particularly for small $|m|$.

\subsection*{Exercise 3: checking a Hartree-Fock code}
\label{ex:11-hfcheck}
\addcontentsline{toc}{subsection}{Exercise 3: checking a Hartree-Fock code}

\begin{enumerate}
  \item[a)] Show that for $N=2$ in the minimal basis the Hartree-Fock energy
  is $2\hbar\omega+\sqrt{\pi/2}\sqrt{\hbar\omega}$, and check it against
  Table~\ref{tab:11-hfconvergence}.
  \item[b)] Explain why the $N=2$ energy does not change from one shell to
  two, nor from three to four.
  \item[c)] What would you conclude if the converged single-particle energies
  within a shell were not degenerate?
\end{enumerate}

\paragraph{Answer to a).}
With only the $(0,0)$ spatial orbital available, both electrons occupy it with
opposite spins and no variational freedom remains: $\bm{C}$ is the identity.
The energy is then the unperturbed $2\hbar\omega$ plus the single direct
matrix element $\bracket{00;00}{\hat{v}}{00;00}$ -- the exchange term vanishes
because the two electrons have opposite spins -- which is
Eq.~(\ref{eq:11-lowestelement}).  At $\hbar\omega=1$ this is
$2+1.25331414=3.25331414$, the first row of the table.

\paragraph{Answer to b).}
The Hartree-Fock orbitals carry good angular momentum, since $\hat{f}$
commutes with $\hat{L}_z$, so the occupied $m=0$ orbital can mix only with
other $m=0$ states.  Shell $N_s=1$ contains none: its states are $(0,\pm1)$.
Adding it therefore changes nothing.  The first new $m=0$ state is $(1,0)$ in
shell $2$, which is why the energy drops there, and the next is $(2,0)$ in
shell $4$, which is the second drop.  Shells $3$ and $5$ contain no $m=0$
states at all.  The same argument explains why the two-electron energy is
converged to six digits long before the six-electron one.

\paragraph{Answer to c).}
Either the calculation has not converged, or the antisymmetrised matrix
elements are wrong in a way that breaks rotational invariance -- the most
common cause being the spin deltas of Eq.~(\ref{eq:11-antisym}), where the
exchange term has $r$ and $s$ interchanged.  A broken degeneracy is a far more
sensitive diagnostic than the total energy, which can be wrong by a little
without looking wrong at all.

\subsection*{Exercise 4: the singles amplitudes}
\label{ex:11-singles}
\addcontentsline{toc}{subsection}{Exercise 4: the singles amplitudes}

\begin{enumerate}
  \item[a)] State precisely what Brillouin's theorem says, and what it does
  not say.
  \item[b)] Show that $t_1=0$ is \emph{not} a solution of the CCSD singles
  equation for an interacting system, using
  Eq.~(\ref{eq:11-singlesdriver}).
  \item[c)] Design a numerical test that would catch a CCSD implementation in
  which the two driver terms have been omitted.
\end{enumerate}

\paragraph{Answer to a).}
Brillouin's theorem says that the Hartree-Fock determinant does not couple to
singly excited determinants through the Hamiltonian,
$\bracket{\Phi_i^a}{\hat{H}}{\Phi_0}=f_{ai}=0$.  That is a statement about a
single matrix element.  It implies that singles do not contribute at first
order in perturbation theory, and hence that MP2 involves only doubles.  It
says nothing about the coupled-cluster amplitudes, which solve a non-linear
system in which the singles are coupled to the doubles.

\paragraph{Answer to b).}
Set $t_1=0$ in the singles equation.  Every term containing $t_1$ vanishes and
so does $f_{ai}$, but the two terms of Eq.~(\ref{eq:11-singlesdriver}) are
built from $t_2$ alone and are generally non-zero.  The residual is therefore
not zero and $t_1=0$ is not a fixed point.  Iterating from $t_1=0$ generates
singles of the order of the doubles times a matrix element over an energy
denominator, which is why they come out small -- about $10^{-2}$ here.

\paragraph{Answer to c).}
Run the code on a two-electron system.  Singles and doubles exhaust the
excitation space for two particles, so CCSD must equal full configuration
interaction exactly.  A solver missing the driver terms will converge to
something close but not equal: the discrepancy in
Table~\ref{tab:11-twoelectron} would be of order $10^{-4}$ rather than
$10^{-13}$, and no amount of tightening the convergence threshold will remove
it.  The same test distinguishes a genuine convergence problem, which
tightening does fix, from a wrong equation, which it does not.

\subsection*{Exercise 5: numerical explorations}
\addcontentsline{toc}{subsection}{Exercise 5: numerical explorations}

\begin{enumerate}
  \item[a)] Using \texttt{quantumdot.py}, reproduce
  Table~\ref{tab:11-twoelectron} and extend it to seven shells.  Fit the
  remaining error to $a/M+b/M^2$ and extrapolate to the complete basis.  How
  close do you get to $3$?
  \item[b)] Compute the twelve-electron dot at $\hbar\omega=1$ in the full
  basis.  How does the fraction $E_{\mathrm{CCSD}}/(E_{\mathrm{HF}}-E_0)$
  compare with the two- and six-electron values, and why?
  \item[c)] Plot the Hartree-Fock and bare oscillator single-particle spectra
  for $N=6$ and comment on Koopmans' theorem.
  \item[d)] For the six-electron dot in the two-shell basis, optimise the
  unitary amplitudes properly rather than borrowing them from CCSD.  How much
  further does the energy fall, and how many function evaluations does it
  take?
  \item[e)] Repeat the Trotter study of Table~\ref{tab:11-trotter} with the
  symmetric second-order splitting of Eq.~(\ref{eq:6-strang}) and confirm that
  the error now falls as $1/n^2$.
  \item[f)] Estimate the number of qubits and Pauli strings needed for the
  six-electron dot in the two-shell basis, using
  Section~\ref{sec:jordanwigner}.  Compare with the pairing model of
  Table~\ref{tab:4-paulicounts}.
\end{enumerate}
